% Generated by Listas/make_statements.py from pratica-lecture-08.md, item 23. Do not edit.
\textbf{Source.} Muñoz-Garcia \& Toro-Gonzalez (2019), \emph{Strategy and Game Theory: Practice Exercises with Answers}, Ch. 11, Exercise 4 ``Labor Market Signaling Game and Equilibrium Refinements'', p.~500 {[}PDF p.~507{]}. Reproduced verbatim (statement only; the book's worked answer follows it). Answer: yes, the book's worked answer follows the exercise.



\begin{quote}

Consider the following labor market signaling model. A worker with privately known productivity \(\theta\) chooses an education level \(e\). Upon observing \(e\), each firm responds with a wage offer that corresponds to the worker's expected productivity (That is, we assume that the labor market is perfectly competitive, so firms pay the value of the expected marginal productivity to every worker. Otherwise, competing firms would have incentives to increase their salaries in order to attract workers with higher productivities.). The worker's payoff from receiving a wage \(w\) and acquiring an education level \(e\), given an innate productivity \(\theta\), is \(u(w, e \mid \theta) = w - \frac{e^3}{\theta}\), i.e., the cost of acquiring education, \(\frac{e^3}{\theta}\), is convex, and decreases in the worker's innate productivity, \(\theta\) (Note that, for simplicity, we assume that the worker's innate productivity, \(\theta\), is unaffected by the education level he acquires.).



\textbf{Part (a)} Assuming that there are two equally likely types \(\theta_1 = 1\) and \(\theta_2 = 2\), characterize the set of all separating equilibria.



\textbf{Part (b)} Which of these equilibria survive the Cho and Kreps' (1987) Intuitive Criterion?

\end{quote}

